% Part: model-theory
% Chapter: interpolation
% Section: introduction

\documentclass[../../../include/open-logic-section]{subfiles}

\begin{document}

% SEGMENT OLP-0199-S01

\olfileid{mod}{int}{int}

\olsection{அறிமுகம்}

இடையீட்டுத் தேற்றம் பின்வரும் முடிவு: $\Entails !A \lif !B$ என்று
கொள்க. அப்போது $\Entails !A \lif !C$ மற்றும்
$\Entails !C \lif !B$ ஆகுமாறு !!a{sentence} $!C$ ஒன்று உள்ளது.
மேலும், $!C$ இலுள்ள ஒவ்வொரு !!{constant}, !!{function},
!!{predicate} உம் ($\eq$ தவிர) $!A$ மற்றும்~$!B$ ஆகிய இரண்டிலும்
இடம்பெறும். !!{sentence} $!C$ ஆனது $!A$ மற்றும்~$!B$ இன்
\emph{இடையீட்டி} எனப்படுகிறது.

இடையீட்டுத் தேற்றம் தன்னளவிலேயே சுவையானது; ஆனால் ஒரு கோட்பாட்டில்
வரையறுக்கத்தக்கமை குறித்த முடிவுகளையும், முரணின்மையுடைய இரு
கோட்பாடுகளை இணைத்தால் எந்நிபந்தனைகளில் முரணின்மையுடைய கோட்பாடு
கிடைக்கும் என்பது குறித்த முடிவுகளையும் நிரூபிக்க அதைப் பயன்படுத்த
முடிவதே அதன் முதன்மைச் சிறப்பு. முதலாவது பெத்தின் வரையறுக்கத்தக்கமைத்
தேற்றம் என்றும், இரண்டாவது ராபின்சனின் கூட்டு முரணின்மைத் தேற்றம்
என்றும் அறியப்படுகின்றன.

\end{document}
